\begin{definition}[精确高表示的数据]
\label{def:stage-high-data}
固定实数 $l_x,l_y>0$ 与 $\ell\ge0$，使得
\begin{equation}
 X=Z^{l_x},\qquad Y=Z^{l_y},\qquad h=1-l_x+\ell>0,
 \qquad
 C(s)=s+\frac{l_x}{2}-1+\frac h6
     =s-\frac56+\frac{l_x}{3}+\frac\ell6.
 \label{eq:stage-mellin-exponent}
\end{equation}
这些实数在目标特征之前固定。对每个本原有限阶目标 $\eta$，固定节~\ref{sec:probe} 的
数据 $S,\xi,W_0,W_1$，它与 $Z$ 无关。精确高表示的数据由以下对象与断言构成。

首先，对每个充分大的实数 $Z$，$\mathscr I_\eta(Z)$ 是一个独立指定的复值
量。在每次应用中，它由该 $Z$ 下方程~\eqref{eq:general-probe} 中完备化和的有限
线性组合给出，允许对完备化指标 $A$ 作指定限制，并对其三个正尺度作指定的
重标度。该有限组合中的项数与系数可依赖于 $Z$。任何对 $\mathscr I_\eta$ 的直接
估计都指这一表达式，而非某个另行定义的高积分。

其次，对每个满足 $(u,S)=1$ 的非零无六次幂因子元素 $u$ 及每个这样的 $Z$，
存在函数 $\mathfrak H_{\eta,u,Z}(s,w,z)$。它在 $\mathcal D_2$ 的每一点以及每个
$\mathcal D_1(\epsilon_0)$ 的每一点的邻域上全纯。对每个 $D>0$ 及包含于这些
区域之一的每个实长方体 $\mathcal K$，存在有限常数 $B_{\eta,\mathcal K,D}$ 与
$J_{\eta,\mathcal K,D}$，使得对 $1\le q_u\le Z^D$、所有虚部以及 $\mathcal K$ 中的
实部一致地有
\begin{equation}
 |\mathfrak H_{\eta,u,Z}(s,w,z)|
 \ll_{\eta,\mathcal K,D}
 Z^{B_{\eta,\mathcal K,D}}
 (1+|\Im s|+|\Im w|+|\Im z|)^{J_{\eta,\mathcal K,D}}.
 \label{eq:stage-all-height-majorant}
\end{equation}
这些常数与指数独立于 $Z,u$ 以及任何后续的分部积分次序。不假设 $\mathfrak H$
或其任何局部因子非零。

最后，对独立给定的 $\mathscr I_\eta(Z)$ 断言如下恒等式，其中和式与积分在所示
直线上绝对收敛：
\begin{equation}
\begin{split}
 \mathscr I_\eta(Z)=\frac1{(2\pi i)^3}
 \int_{(3)}\int_{(3)}\int_{(2)}
 &\mathcal W(X,Y,Z;s,w,z)
 \sum_u^{(6)}q_u^{-z}\overline{\xi(u)}\\
 &\cdot
 \frac{\zeta_F^S(6z)L^S(w,\chi_\bullet(u))}
      {L^S(s,\eta\overline{\chi_\bullet(u)})}
 \mathfrak H_{\eta,u,Z}(s,w,z)\,dz\,dw\,ds,
\end{split}
\label{eq:stage-high-identity}
\end{equation}
其中和式具有方程~\eqref{eq:probe-poisson-identity} 的物理限制，并且
\[
 \mathcal W(X,Y,Z;s,w,z)
 =X^{1/2-z}Z^{s+z-1}Y^{w-1}
       e^{(s+z-1)^2}M(z)\widehat W_1(w).
\]
因此方程~\eqref{eq:stage-high-identity} 是对物理表达式待证明的一个假设，而非
其定义。
\end{definition}

对基础探针，物理表达式为
$I_\eta(Z^{l_x},Z^{l_y},Z)$、$\ell=0$ 以及
$\mathfrak H_{\eta,u,Z}=\mathcal H_{\eta,u}$。
方程~\eqref{eq:probe-poisson-identity} 与~\eqref{eq:probe-high}
证明所需的恒等式，而引理~\ref{lem:local-euler} 给出全纯性以及
方程~\eqref{eq:stage-all-height-majorant}。更一般地，定义中的函数
$\mathfrak H$ 始终指
\emph{所示标量商被除去之后的完整修正}。
若该修正是局部因子乘积之和，则这些乘积本身
在所述区域上必须全纯。用个别局部因子在其零点处
作商不能确立这一条件。仅用于估计被保留围道的
分解将在下文单独陈述。

\subsection{竖直线上的外部尾部与水平衔接}

我们三个 Mellin 检验函数的乘积在三个独立的高度方向上衰减。下述估计把这一
衰减转化为大长方体之外以及围道矩形水平边上的界。

\begin{lemma}[外部积分尾部与迹尾部]
\label{lem:external-contour-tails}
设 $m\ge2$，设 $\mathcal K$ 为实围道参数的紧集，且对 $\boldsymbol r\in\mathcal K$ 与
$Z\ge2$，设 $K_{\boldsymbol r,Z}$ 为 $\mathbb R^m$ 上的非负 Borel 可测核。假设它在
$(\boldsymbol r,\boldsymbol t)$ 中联合 Borel 可测，且对每个整数 $L\ge0$，
\[
 K_{\boldsymbol r,Z}(\boldsymbol t)
 \ll_L(1+|\boldsymbol t|)^{-L}
\]
对 $\boldsymbol r,Z,\boldsymbol t$ 一致成立。设
$\mathcal D\subseteq\mathbb R^m$ 为 Borel 可测，并假设 Borel
可测因子 $F_Z$ 在 $\mathcal D$ 上满足
\begin{equation}
 |F_Z(\boldsymbol t)|\le C Z^B(1+|\boldsymbol t|)^J,
 \label{eq:external-arithmetic-majorant}
\end{equation}
其中 $B,J,C$ 固定，独立于 $Z$ 与下文整数 $N$ 之下。
用 $\widehat{\boldsymbol t}_j$ 表示略去 $\boldsymbol t$ 的
第 $j$ 个坐标所得的向量。对 $T\ge1$ 与
$\boldsymbol r\in\mathcal K$ 一致地，
\begin{align}
 \int_{\mathcal D\cap\{\max_j|t_j|>T\}}
 K_{\boldsymbol r,Z}(\boldsymbol t)|F_Z(\boldsymbol t)|\,d\boldsymbol t
 &\ll_N Z^B T^{-N},
 \label{eq:external-integrated-tail}\\
 \int_{\mathcal D\cap\{t_j=T\}}
 K_{\boldsymbol r,Z}(\boldsymbol t)|F_Z(\boldsymbol t)|\,
 d\widehat{\boldsymbol t}_j
 &\ll_N Z^B(1+|T|)^{-N}.
 \label{eq:external-trace-tail}
\end{align}
第二个断言对 $T$ 的任一符号与每个 $j$ 成立，在 $t_j=T$ 上使用坐标
Lebesgue 测度。当假设在该处一致成立，且核、区域指示函数与
算术因子在该实参数与其余坐标中联合 Borel 可测时，两个断言在关于
有界长度实围道区间积分之后仍成立。
\end{lemma}

\begin{proof}
把核的界乘以
方程~\eqref{eq:external-arithmetic-majorant}。在长方体之外对
$(1+|\boldsymbol t|)^{J-L}$ 积分，当 $L>N+J+m$ 时证明
方程~\eqref{eq:external-integrated-tail}。
在 $t_j=T$ 上，对其余 $m-1$ 个坐标积分给出
\[
 \int_{\mathbb R^{m-1}}
 (1+|T|+|\widehat{\boldsymbol t}_j|)^{J-L}
 d\widehat{\boldsymbol t}_j
 \ll(1+|T|)^{J-L+m-1}.
\]
取 $L>N+J+m-1$ 证明方程~\eqref{eq:external-trace-tail}。
限制到 $\mathcal D$ 会减小这些正积分，而有界
实区间只贡献其长度。
\end{proof}

对我们的三重积分，取
\[
 (y_1,y_2,y_3)=(\Im(s+z),\Im z,\Im w),\qquad
 K_{\boldsymbol r,Z}(\boldsymbol t)
 =e^{-y_1^2}|M(r_z+iy_2)|\,|\widehat W_1(r_w+iy_3)|.
\]
高度变量的变换的行列式为 1。在固定的实长方体上，其中 $r_z$ 属于
$(0,\infty)$ 的紧子区间，节~\ref{sec:poisson-representation} 的 Mellin 估计给出 $M$ 与
$\widehat W_1$ 的任意多项式衰减。Gaussian 具有同样的性质。因此它们的乘积
对每个 $L$ 为 $O_L((1+|\boldsymbol y|)^{-L})$，从而也是
$O_L((1+|\boldsymbol t|)^{-L})$。被略去的因子
$e^{(\Re(s+z)-1)^2}$ 在该实长方体上有界。

取 $w$ 的留数之后只剩下 $s,z$。同样的证明适用于核
\begin{equation}
 K^{(2)}_{\boldsymbol r,Z}(\boldsymbol t)
 =e^{-y_1^2}|M(r_z+iy_2)|,
 \qquad(y_1,y_2)=(\Im(s+z),\Im z).
 \label{eq:external-two-dimensional-kernel}
\end{equation}
这提供了后续 $z$-衔接所需的迹界。

同样的核检验处理额外的外部分离变量。
对 $m\ge3$，设 $\boldsymbol y=A\boldsymbol t$，其中 $A\in\mathrm{GL}_m(\mathbb R)$ 固定，
并把三因子核乘以 Borel 密度 $|v_{\boldsymbol r,Z}(y_4,\ldots,y_m)|$，
它在 $(\boldsymbol r,\boldsymbol v)$ 中联合 Borel，满足
\begin{equation}
 \sup_{\boldsymbol r,Z,\boldsymbol v}
 (1+|\boldsymbol v|)^L|v_{\boldsymbol r,Z}(\boldsymbol v)|<\infty
 \quad\text{对每个整数 }L\ge0.
 \label{eq:external-density-contract}
\end{equation}
对 $m=3$，该密度在零维空间上为 1。由于 $A$ 及其逆固定，乘积
在所有 $m$ 个坐标中仍快速衰减。对二因子核，相同的论证在
$\mathbb R^{m-2}$ 上使用密度并要求 $m\ge2$，当 $m=2$ 时密度为 1。
因此两种乘积核形式都具有上述积分界与迹界。

由对数 Fourier 反演分离的一致环形剖面通过反复
分部积分满足方程~\eqref{eq:external-density-contract}，如引理~\ref{lem:smooth-calculus} 中
那样。首先把每个纯范数扭转翻译为其 Mellin 变元。额外的坐标高度
随后向三维变量变换附加一个具有固定逆的块三角矩阵。

仅通过固定的加权 $L^1$ 范数而知的内部系数测度并不
因此被逐点密度假设或其迹结论所覆盖。这样的测度在因子 $F_Z$ 内部
仍以其已经固定的矩阶被积分。任何使用迹估计的额外外部
坐标必须单独满足方程~\eqref{eq:external-density-contract}。

引理中的区域 $\mathcal D$ 很重要。假设某个倒数仅在其变元
$\gamma+\sum_j a_jt_j$ 的虚部位于带缓冲区间内时才有界。具有
$a_j\ne0$ 的坐标只能在保持这一条件的范围内积分；引理不会把
倒数的界扩展到该范围之外。当其他因子在该处满足
方程~\eqref{eq:external-arithmetic-majorant} 时，具有 $a_j=0$ 的坐标可以被扩展。
扩展轴上的所有因子必须使用其全局或绝对估计，而非仅在保留
区间上成立的估计。特别地，仅凭积分尾部不足以证明水平衔接的
合理性：该衔接使用方程~\eqref{eq:external-trace-tail}。

在本引理中增大 $N$ 只会增大外部光滑检验函数的衰减阶。
它不对 $F_Z$ 求导。因此，当 $N$ 稍后被选取时，$F_Z$ 中
固定的矩剖面阶或固定的算术高度指数保持不变。这与
引理~\ref{lem:smooth-calculus} 中所作的阶的区分相同。

\subsection{移动一个固定分箱}

我们首先移动整个行积分。行包络只会在这一步之后，
在被保留的围道上应用。

\begin{lemma}[固定分箱的围道变换]
\label{lem:fixed-bin-contour}
假设定义~\ref{def:stage-high-data}。固定
$\sigma_0\in[7/8,1)$ 且 $\beta_*>\sigma_0$，令 $z_0=17/50$，并设
$U=Z^d$，其中 $0<d_{\min}\le d\le d_{\max}<\infty$。
设 $\mathcal B$ 为引理~\ref{lem:buffered-bins} 的某个分箱 $(i,a)$ 中满足 $q_u\asymp U$
的被保留物理行的有限集，它在任何
Mellin 变量被移动之前固定。该环形区域中的比较常数
固定。取 $0<e<10^{-3}$ 与 $T_1=Z^\tau>2$，其中
$0<\tau\le d_{\min}/100$。实范围与 $e$ 独立于
目标。

固定正常数 $c_s,c_w,c_z\le1/2$，独立于 $Z$ 与
各行。保留原始高度
$|\Im s|\le c_sT_1$、$|\Im w|\le c_wT_1$ 与
$|\Im z|\le c_zT_1$。$\mathcal B$ 在方程~\eqref{eq:stage-high-identity} 中
起始直线上的贡献等于其在这些保留线段上的积分
\begin{equation}
 \Re s=a+16e,\qquad \Re w=1-a-6e,\qquad \Re z=z_0,
 \label{eq:stage-central-contours}
\end{equation}
误差至多 $O_{\eta,N}(Z^{B_\eta}T_1^{-N})$，对每个固定的 $N$。有限
指数 $B_\eta$ 在 $N$ 之前固定。空的行集贡献为零。
该变换使用完整的全纯修正；不需要对该修正
作分解，也不需要依赖于围道点的细分。
\end{lemma}

\begin{proof}

节~\ref{sec:detector} 给出 $a\le\beta_*$，包括地板分箱。
在方程~\eqref{eq:stage-high-identity} 的绝对直线上隔离出关于 $\mathcal B$
的有限和。把
$z$ 从 $2$ 移动到 $z_0$，保持 $w=3$，并把 $s$ 从 $3$ 移动到
$\beta_*+20e$。这些移动对某个固定的正 $\epsilon_0$ 保持在 $\mathcal D_1(\epsilon_0)$
内，两个分子变元保持在其极点右侧，且倒数保持在
$\Re s>\beta_*$。其全局界由引理~\ref{lem:logarithmic-control}
与~\ref{lem:deleted-euler-factors} 得出，使用节~\ref{sec:detector} 中的
导子与删除根式界。该修正服从
方程~\eqref{eq:stage-all-height-majorant}。因此 $T\to\infty$ 的
引理~\ref{lem:external-contour-tails} 证明这些移动是合理的。

现在把 $w$ 从 $3$ 移动到 $1-a-6e$，而 $s$ 保持在该全局直线上。
在整个移动过程中，
\[
 \Re(s+w)\ge1+(\beta_*-a)+14e\ge1+14e.
\]
每个被保留行的分子特征是非主的，因此其
$L$ 函数是整函数。又有 $\Re w\ge-6e>-1/100$ 与
$\Re z=z_0$。于是该修正在
$\mathcal D_1(\epsilon_0)$ 中保持全纯，例如取 $\epsilon_0=10e$，且不跨
越任何极点。在 $w$-衔接上倒数仍是全局的。分子的
全局上带形域界、删除因子界以及
方程~\eqref{eq:stage-all-height-majorant} 给出在该处所积分坐标轴上的形如
方程~\eqref{eq:external-arithmetic-majorant} 的优函数。迹估计使这些衔接趋于零。

在把 $s$ 移动得更远之前，把三个原始 Mellin 高度限制到所述长方体。在被丢弃的
部分上保持 $s$ 在 $\Re s=\beta_*+20e$。在那里，全局倒数界、
全局分子界、删除因子以及全高度
修正界给出 $Z^{B_\eta}$ 乘以高度的固定多项式：
行范围由 $q_u\ll Z^{d_{\max}}$ 界定，而行数
为 $O(U)$。引理~\ref{lem:external-contour-tails} 中的积分尾部给出
$O_{\eta,N}(Z^{B_\eta}T_1^{-N})$。

只把被保留的 $s$ 线段移动到 $a+16e$。在这个矩形上，
\[
 \Re(s+w)\ge1+10e,
\]
且定义 $\mathcal D_1(10e)$ 的其他不等式仍然成立。
倒数变元的实部至少为 $a+16e>a+6e$。
所述高度长方体对充分大的 $Z$ 把两个标量变元
严格置于 $|\Im v|\le(3i+2)T_1$ 内部。
因此引理~\ref{lem:buffered-bins} 在整个
被保留矩形上排除其零点。由于 $6z_0>1$，标量 zeta 因子在此移动上
没有极点。于是不跨越任何极点。

对 $s$-衔接，其虚坐标固定为 $T_1$ 的常数倍。
方程~\eqref{eq:stage-high-identity} 中的标量分母
只依赖于 $s$，因此当
$\Im z$ 与 $\Im w$ 的积分扩展到其整个坐标轴时，其带缓冲的估计仍然有效。
在这些扩展轴上使用全局分子估计与
方程~\eqref{eq:stage-all-height-majorant}。迹估计于是给出
$O_{\eta,N}(Z^{B_\eta}T_1^{-N})$。本论证中没有任何 $w$- 或 $z$-衔接
扩展一个已被移动到仅是带缓冲区域的 $s$-轴。

\end{proof}

精确高表示是一个三重积分，带有一个关于其
完整修正的优函数。因此我们可以在
引入任何辅助分离变量之前执行这一变换。若某个中心估计稍后
使用这类变量，则其保留区域与丢弃部分必须
在该估计中得到证明；它们不是刚刚证明的围道
恒等式的新坐标。

\subsection{应用行包络}

对基础修正，引理~\ref{lem:local-euler} 验证
引理~\ref{lem:fixed-bin-contour} 的假设。事实上
$|\mathcal H_{\eta,u}|\ll_\epsilon q_u^\epsilon$ 在两个 Euler 区域中对所有三个
高度一致成立；对 $q_u\le Z^D$，取 $\epsilon=1$ 给出
方程~\eqref{eq:stage-all-height-majorant}，其中 $B=D$ 且高度阶
为零。我们现在估计被保留的积分，在本小节中始终使用
$l_x=l_y=h=1/2$ 与 $\ell=0$。

固定
\[
 d_{\min}=\frac1{63},\qquad \zeta=\frac1{1000},\qquad
 d_{\max}=h+\zeta=\frac{501}{1000}.
\]
对该范围内的一个二元组 $U=Z^d$，在移动
任何 Mellin 变量之前固定一个动态分箱 $(i,a)$，并令 $\delta=2a-1$。
命题~\ref{prop:sextic-row-count} 在请求损失
$\epsilon_r>0$ 下给出其基数至多
$U^{R(\delta)+\epsilon_r}(1+T_1)^{A_\eta}$，相差一个固定常数。
对地板而言这是直接计数 $O_\eta(U)$，因为
$R(1/50)=1$；它不使用见证。

分子 $L^S(w,\chi_\bullet(u))$ 正是节~\ref{sec:detector} 中
$\nu=1$ 的原始零扩展表示。
在被保留直线 $\Re w=1-a-6e$ 上，引理~\ref{lem:buffered-bins} 对任意
$\epsilon_n>0$ 给出界
$U^{\delta/2+12e+\epsilon_n}(3+T_1)^C$。
这在地板分箱中也成立。中心点位于
$\mathcal D_1(10e)$ 内，因此引理~\ref{lem:local-euler} 中更强的局部估计给出
该处 $|\mathcal H_{\eta,u}|\ll_{e,\epsilon_H}U^{\epsilon_H}$，对任意
$\epsilon_H>0$。于是
\[
 \sum_{u\in\mathcal B}
 |L^S(w,\chi_\bullet(u))\mathcal H_{\eta,u}(s,w,z)|
 \ll_\eta
 U^{R(\delta)+\delta/2+12e+\epsilon_r+\epsilon_n+\epsilon_H}
 (1+T_1)^{A'_\eta}.
\]
这里 $A'_\eta$ 有限，并在最终外部尾部阶之前固定。

令 $\varepsilon_c=12e+\epsilon_r+\epsilon_n+\epsilon_H$。
带缓冲的倒数对任意
$\varepsilon_d>0$ 耗费 $U^{\varepsilon_d}$，且 $\zeta_F^S(6z_0)$ 绝对有界。
这些检验函数在上述三个独立高度
坐标中具有有界的联合 $L^1$ 范数。外部幂，包括 $q_u^{-z_0}$，为
\[
 Z^{a/2-3/4+z_0/2+13e}U^{-z_0}.
\]
把这些界相乘，对每个 $\varepsilon_p>0$ 给出至多的被保留
分箱贡献
\begin{equation}
 \ll_\eta Z^{1/4+E_{\mathrm I}(d)+13e
       +(1+d)\varepsilon_c+d\varepsilon_d+\varepsilon_p}
 (1+T_1)^{A'_\eta}.
 \label{eq:balanced-central-contribution}
\end{equation}
这里我们无害地
把这界扩大了 $Z^{\varepsilon_c+\varepsilon_p}$，使得各项损失使用与
后续一般估计相同的形式，并且
\[
 E_{\mathrm I}(d)=-\frac34+\frac{\delta+R(\delta)}2
 +\left(d-\frac12\right)
       \left(R(\delta)+\frac\delta2-\frac{17}{50}\right).
\]
在此应用中没有额外的 Mellin 积分。其标量
带缓冲变元具有高度 $\Im s$ 与 $\Im w$。行计数中的见证
频率与剖面选择保持在该探测器已有的
容许范围之内，并且不进入任何一个标量变元。

在 $d=h=1/2$ 处，行包络的三片分别给出
\[
 E_{\mathrm I}(h)=
 \begin{cases}
  -1/4+\delta/2,&0\le\delta\le1/3,\\
  -1/12,&1/3\le\delta\le2/3,\\
  -(1-\delta)/4,&2/3\le\delta\le1.
 \end{cases}
\]
前两片至多为 $-1/12$。在第三片中，节~\ref{sec:detector} 的精确分箱
上界为
\[
 \delta\le2\beta_*-1=\frac56+2\Delta_1.
\]
由此在每一片中都有
\[
 E_{\mathrm I}(h)-\Delta_1
 \le-\frac1{24}-\frac{\Delta_1}{2}.
\]
这就是此处所需的与 $C_{\mathrm I}(\beta_*)$ 的比较；
$E_{\mathrm I}(h)$ 本身不必为负。
频率斜率为三种形式
\[
 R(\delta)+\frac\delta2-\frac{17}{50}
 =\begin{cases}
  33/50+\delta/2,&0\le\delta\le1/3,\\
  149/150-\delta/2,&1/3\le\delta\le2/3,\\
  33/50,&2/3\le\delta\le1.
 \end{cases}
\]
因此它位于 $[33/50,62/75]$ 内。正性控制所有 $d\le h$，
而到 $d\le h+\zeta$ 的延拓最多耗费 $(62/75)\zeta$。
于是，在可调损失之前，每个中心二元组至少有节省
\[
 m_c=\frac1{24}-\frac{62}{75}\zeta=\frac{1021}{25000}
\]
相对于 $C_{\mathrm I}(\beta_*)$。

\subsection{把修正分成若干片段的中心估计}

前面的计算把一个总体的绝对行界乘以
外部 Mellin 幂。我们记录当完整修正由
若干片段估计、且可用界在行集之间变化时它的形式。
围道对固定分箱已经得到证明；这些集合
只用于估计其被保留的被积函数。

\begin{lemma}[一个被保留积分及其指数]
\label{lem:bin-contour-accounting}
假设定义~\ref{def:stage-high-data} 的数据。固定
$\sigma_0\in[7/8,1)$ 且 $\beta_*>\sigma_0$，并令 $z_0=17/50$。
设 $0<d_{\min}<d_{\max}<\infty$，$U=Z^d$ 且
$d_{\min}\le d\le d_{\max}$，并设 $\mathcal B$ 为
引理~\ref{lem:buffered-bins} 的某个被保留分箱 $(i,a)$ 中满足 $q_u\asymp U$ 的
物理行 $u$ 的有限集。$q_u\asymp U$
中的比较常数固定。令 $\delta=2a-1$，取如该引理中的 $0<e<10^{-3}$，
并设 $T_1=Z^\tau>2$，其中 $0<\tau\le d_{\min}/100$。
实范围与 $e$ 的选取独立于目标。

分箱 $\mathcal B$ 在任何 Mellin 变量被移动之前固定。在
方程~\eqref{eq:stage-central-contours} 的被保留中心围道上，假设存在分解
\[
 \mathfrak H_{\eta,u,Z}(s,w,z)
       =\sum_{\lambda\in\Lambda}
             \mathfrak H^{(\lambda)}_{\eta,u,Z}(s,w,z),
\]
其中 $\Lambda$ 有限、非空且在 $Z$ 之前固定。这一等式
只在被保留围道上要求；各求和项不必在远离
这些围道处全纯。对每个 $\lambda$ 与每个被保留的 $(s,w,z)$，假设
$\mathcal B$ 被划分成至多 $C(1+\log Z)^D$ 个集合
$\mathcal B_{\lambda,\nu}(s,w,z)$，其中 $C,D$ 固定。对每个集合在
每个被保留点，设 $R,g$ 为某个固定有界范围内的实数。假设，对固定的
$\varepsilon_c\ge0$ 与有限的 $A_\eta\ge0$，
\begin{equation}
 \sum_{u\in\mathcal B_{\lambda,\nu}(s,w,z)}
  \left|L^S(w,\chi_\bullet(u))
       \mathfrak H^{(\lambda)}_{\eta,u,Z}(s,w,z)\right|
 \ll_\eta U^{R+\delta/2+\varepsilon_c}
 Z^{\ell(z_0-1/2)+g+\varepsilon_c}(1+T_1)^{A_\eta}.
 \label{eq:stage-central-bound}
\end{equation}
该界在被保留点、行集与移动标签中一致。
数 $\varepsilon_c$ 与 $R,g$ 的有界范围在
目标之前固定；$R,g$ 可依赖于 $d,a$、逐点集合与全局数
$\beta_*$，但不以其他方式依赖于 $\eta$。设 $\mathcal P_Z$ 为
该分解中在任何被保留点出现的所有
数对 $(R,g)$ 的集合。

这里“被保留”具有如下精确的高度含义。一次列出
在验证
方程~\eqref{eq:stage-central-bound} 时进入某个带缓冲的 $L$-值、
倒数或对数导数的所有外部坐标 $v_1,\ldots,v_m$，并
包含所有三个原始
Mellin 高度，当它们不出现在某个变元中时系数取零。经过
纯扭转平移之后，每个这样的变元
具有虚部
$\gamma_j+\sum_{k=1}^m a_{jk}v_k$，其中有限矩阵 $(a_{jk})$
固定且 $|\gamma_j|\le(3i+1)T_1$。限制
\begin{equation}
 |v_k|\le c_kT_1,\qquad
 c_k=\frac1{2m(1+\max_j|a_{jk}|)}.
 \label{eq:stage-height-allocation}
\end{equation}
系数为零的坐标可以按同样的规则限制。
于是每个新增高度的总和至多为 $T_1/2$。
固定的有界扩大通过提高 $Z$ 的下阈值来纳入。
若使用额外的 Mellin 积分来确立
方程~\eqref{eq:stage-central-bound}，则该不等式在那些积分之后
对其整个左侧被要求成立。任何被丢弃的部分都必须
留在其全局或绝对直线上，由
引理~\ref{lem:external-contour-tails} 或
引理~\ref{lem:smooth-calculus} 界定，并在
假设被断言之前包含在所示界中。容许量 $T_1/2$ 在
接连的估计中不会被更新。

与数对 $(R,g)$ 相关联的公共理想指数为
\begin{equation}
 E_{\sigma_0}(d;R,g)=a-\sigma_0+h(z_0-1/6)-a l_y-\ell/2+g
                  +d(R+\delta/2-z_0).
 \label{eq:stage-high-exponent}
\end{equation}
若 $\mathcal B$ 为空，则其贡献为零。否则令
\[
 E^{\max}_{\sigma_0,Z}(d)
   =\sup_{(R,g)\in\mathcal P_Z}E_{\sigma_0}(d;R,g).
\]
对每个 $\varepsilon_d,\varepsilon_p>0$，$\mathcal B$ 在中心围道上的
总贡献对充分大的 $Z$ 为
\begin{equation}
 \begin{split}
 \ll_\eta{}& Z^{C(\sigma_0)+E^{\max}_{\sigma_0,Z}(d)+(16-6l_y)e}\\
 &\cdot Z^{(1+d)\varepsilon_c+d\varepsilon_d+\varepsilon_p}
       (1+T_1)^{A_\eta}.
 \end{split}
 \label{eq:stage-central-contribution}
\end{equation}
上确界有限，因为这些数对位于一个固定有界范围内。
幂 $\varepsilon_p$ 吸收逐点对数重数。
不声明关于逐点片段或行集的单独积分。
此围道移动中所有被丢弃部分与水平衔接对每个固定 $N$ 都是
$O_{\eta,N}(Z^{B_\eta}T_1^{-N})$，其中 $B_\eta$
有限且在 $N$ 之前固定。
\end{lemma}

\begin{proof}
把引理~\ref{lem:fixed-bin-contour} 应用于原始三重积分，
对其三个长方体常数使用方程~\eqref{eq:stage-height-allocation} 中
相应的 $c_k$。这些常数为正且至多为 $1/2$。
因此完整修正与固定分箱在任何逐点片段
被引入之前给出被保留积分与所断言的误差。任何用于中心界的
辅助积分都受陈述中额外假设的约束。

现在使用中心分解并构造
逐点集合 $\mathcal B_{\lambda,\nu}$。在每个被保留点对
原始行和施加三角不等式，然后对
该点存在的集合应用
方程~\eqref{eq:stage-central-bound}。
逐点至多有 $(1+\log Z)^D$ 的固定倍数个这样的集合。
把它们的每个指数用 $\mathcal P_Z$ 上的上确界来界定，并
只对原始全纯行和积分。于是既不需要个别逐点细分的
可测性，也不需要横跨围道点的一个公共标签集。在
被保留围道上，
引理~\ref{lem:buffered-bins} 把倒数界定为
$O_{\eta,e,\varepsilon_d}(U^{\varepsilon_d})$。这些检验函数具有
一致有界的联合 $L^1$ 范数，因为变换
$(\Im s,\Im z,\Im w)\mapsto(\Im(s+z),\Im z,\Im w)$ 可逆。
标量因子 $\zeta_F^S(6z_0)$ 绝对有界。剩下要
计算 $Z$ 的实幂。

